\section{Comparison theorems for \'etale coverings}
\label{XII.5}
\setcounter{subsection}{-1}
\subsection{}
\label{XII.5.0}
Let us make precise the notion of finite covering of an analytic space.
If $\othercal{X}$ is an analytic space, one says that an analytic space
$\othercal{X}'$ finite over~$\othercal{X}$ is a finite covering of
$\othercal{X}$ if every irreducible component of~$\othercal{X}'$ dominates
an irreducible component of~$\othercal{X}$.

\begin{theorem}[``Riemann existence theorem'']
\label{XII.5.1}
Let $X$ be a $\CC$-scheme locally of finite type, $X^{\an}$
the analytic space associated with $X$. The functor $\Psi$ which,
to every finite \'etale covering $X'$ of~$X$, associates
$X^{\prime\an}$ is an equivalence
% original p. 333
from the category of finite \'etale coverings of~$X$ to
the category of finite \'etale coverings of~$X^{\an}$.
\end{theorem}

1) \emph{The functor $\Psi$ is fully faithful.} Let $X'$
and $X''$ be two finite \'etale coverings of~$X$, and let us prove that
the canonical map
\begin{equation*}
\label{eq:XII.5.*}
\tag{$*$} { \Hom_{X}(X',X'') \to
\Hom_{X^{\an}}(X^{\prime\an},X^{\prime\prime\an}) }
\end{equation*}
is bijective. We may assume $X'$ connected. To give an
$X$-morphism from~$X'$ to $X''$ is equivalent to giving a
connected component $X_{i}$ of~$X'\times_{X}X''$ such that the morphism
$X_{i}\to X'$ induced by the first projection is an
isomorphism. Since the connected components of~$X'\times_{X}X''$
correspond bijectively to the connected components of
$X^{\prime\an}\times_{X^{\an}}X^{\prime\prime\an}$~\eqref{XII.2.6} and since a
morphism $X_{i}\to X'$ is an isomorphism if and only if the same
holds for~$X_{i}^{\an}\to X^{\prime\an}$, this proves the
bijectivity of~\eqref{eq:XII.5.*}.

2) \emph{The functor $\Psi$ is essentially surjective.} Let
$\othercal{X}'$ be a finite \'etale covering of~$X^{\an}$ and let us prove
that there exists an \'etale covering $X'$ of~$X$ such that we have
an isomorphism $X^{\prime\an}\isomto\othercal{X}'$. Taking~1) into account, the
question is local on~$X$, and we may therefore assume $X$ affine.

a) Reduction to the case where $X$ is normal. We may assume $X$
reduced. Indeed, suppose the theorem proved
for $X_{\red}$. The functor which, to a finite \'etale
covering $X'$ of~$X$, assigns the finite \'etale covering
$X^{\prime\an}_{\red}$ of~$X_{\red}^{\an}$ is then an equivalence.
Since it is obtained by composing $\Psi$ with the functor $\Theta$ which,
to a finite \'etale covering of~$X^{\an}$, associates its inverse
image on~$X_{\red}^{\an}$, and since $\Theta$ is fully
faithful, this shows that $\Psi$ is an equivalence of
categories.

We
% original p. 334
may assume $X$ normal. Indeed, let $\widetilde{X}$ be the
normalization of~$X$, $p\colon\widetilde{X}\to X$ the canonical
morphism. Since $p$ is finite, $p$ is an effective descent
morphism for the category of \'etale
coverings~(IX~\Ref{IX.4.7}). The theorem being
assumed proved for $\widetilde{X}$, if we set
$\widetilde{\othercal{X}'}=\othercal{X}'\times_{X^{\an}}\widetilde{X}^{\an}$,
there exist an \'etale covering $\widetilde{X}'$ of
$\widetilde{X}$ and an isomorphism $\widetilde{X}'{}^{\an}\simeq
\widetilde{\othercal{X}'}$. It then follows from~1) that the natural
descent datum that we have on~$\widetilde{\othercal{X}'}$
lifts to a descent datum on~$\widetilde{X}'$
relative to $\widetilde{X}\to X$; this proves the existence of an
\'etale covering $X'$ of~$X$ such that we have an isomorphism
$i\colon X^{\prime\an}\times_{X^{\an}}\widetilde{X}^{\an} \simeq
\widetilde{\othercal{X}'}$ whose inverse images under the two
projections of
$\widetilde{X}^{\an}\times_{X^{\an}}\widetilde{X}^{\an}$ are the
same. By~IX~\Ref{IX.3.2}, whose proof is
valid in the analytic case, the morphism $\widetilde{X}^{\an}\to
X^{\an}$ is a descent morphism for the category of
\'etale coverings, and consequently $i$ comes from an
isomorphism $X^{\prime\an}\simeq
\othercal{X}'$.

b) Reduction to the case where $X$ is regular. Let $U$ be
the open set of regular points of~$X$, $i\colon U\to X$,
$i^{\an}\colon U^{\an}\to X^{\an}$ the canonical morphisms; since
$X$ is normal, we have $\codim(X-U,X)\geq 2$. Suppose that there exists an
\'etale covering $U'$ of~$U$ such that we have
$U^{\prime\an}\simeq \othercal{X}'|U^{\an}$, and let us show that then $U'$
extends to an \'etale covering $X'$ of~$X$ such that we have
$X^{\prime\an}\simeq \othercal{X}'$. It suffices to see that $U'$ extends
to an \'etale covering $X'$ of~$X$; indeed we shall then have an
isomorphism $X^{\prime\an}|U^{\an}\simeq \othercal{X}'|U^{\an}$; but, if
$\cal{F}$ and $\cal{G}$ are the coherent sheaves of algebras
over~$\cal{O}_{X^{\an}}$ defining $\othercal{X}'$ and
$X^{\prime\an}$ respectively, the fact that $X$ is normal and that we have
$\codim(X-U,X)\geq 2$ implies that the canonical morphisms
$$
\cal{F}\to i_*^{\an}(\cal{F}|U^{\an}) \qquad
\cal{G}\to i_*^{\an}(\cal{G}|U^{\an})
$$
are
% original p. 335
isomorphisms \cite[\no 3~prop\ptbl 4]{XII.11}. It follows that $\cal{F}$
and $\cal{G}$, hence also $X^{\prime\an}$ and $\othercal{X}'$, are isomorphic.

Let $\varphi\colon X^{\an}\to X$ be the canonical morphism. Since the
problem of extending $U'$ to $X$ is local on~$X$, it suffices
to prove that, for every point $y$ of~$X^{\an}- U^{\an}$, the \'etale covering
$U'_{\varphi(y)}=U'\times_{X}\Spec\cal{O}_{X,\varphi(y)}$ of
$U_{\varphi(y)}=U\times_{X}\Spec\cal{O}_{X,\varphi(y)}$ extends
to $\Spec\cal{O}_{X,\varphi(y)}$. Let $H$ be the coherent
$\cal{O}_{U}$-Algebra defining $U'$. The canonical
morphism
$$
\alpha\colon (i_*H)^{\an} \to i_*^{\an}(H^{\an})
=\cal{F}
$$
defines a morphism of sheaves of modules on~$\Spec\cal{O}_{X^{\an},y}$:
$$
\alpha_{y}\colon (i_*H)^{\an}_{y} \to \cal{F}_{y}\quoi,
$$
whose restriction to $U_{y}=U_{\varphi(y)}
\times_{\Spec\cal{O}_{X,\varphi(y)}}\Spec\cal{O}_{X^{\an},y}$ is an
isomorphism. But this proves that $H|U_{y}$ is trivial, hence that
$U'_{\varphi(y)}$ extends to $\Spec\cal{O}_{X,\varphi(y)}$.

c) Case where $X$ is affine regular. Let
$$
j\colon X\to P
$$
be a compactification of~$X$, where $P$ is a projective
$\CC$-scheme and $j$ a dominant open immersion. Thanks to the
resolution of singularities theorem~\cite{XII.8}, one can
find a regular scheme $R$ and a projective morphism
$r\colon R\to P$ such that $r$ induces an isomorphism $r^{-1}(X)\simeq
X$ and $r^{-1}(X)$ is the complement in $R$ of a
divisor with normal crossings. Let
% original p. 336
$$
k\colon X\to R
$$
be the canonical immersion. We are going to show that there exists a normal
finite covering~\eqref{XII.5.0} $\othercal{R}'$ of~$R^{\an}$ which extends the
\'etale covering $X^{\prime\an}$. By
Proposition~\Ref{XII.5.3} below, such a covering is
unique; the problem of extending $X^{\prime\an}$ is therefore local on~$R^{\an}$ in a neighborhood of~$R^{\an}- X^{\an}$. Now each
point of~$R^{\an}- X^{\an}$ has an open neighborhood
$\othercal{V}$ isomorphic to a ball of an affine space
$\othercal{E}^{n}$, such that $\othercal{V}- \othercal{V}\cap X^{\an}$
is defined by the vanishing of the first $p$ coordinate
functions $z_{1},\dots,z_{p}$, with $0\leq p\leq n$. The fundamental
group of~$\othercal{U}=\othercal{V}\cap X^{\an}$ is isomorphic to
$\ZZ^{p}$, and every \'etale covering of~$\othercal{U}$ is a quotient
of a covering of the form
$$
\othercal{U}''= \othercal{U}[T_{1},\dots,T_{p}]/
(T_{1}^{n_{1}}-z_{1},\dots,T_{p}^{n_{p}}-z_{p})\quoi,
$$
where the $n_{i}$ are integers $>0$, by a subgroup $H$ of the
Galois group $\ZZ/n_{1}\ZZ\times\cdots\times\ZZ/n_{p}\ZZ$ of
$\othercal{U}''$. Now $\othercal{U}''$ extends to the regular
covering
$$
\othercal{V}''= \othercal{V}[T_{1},\dots,T_{p}]/
(T_{1}^{n_{1}}-z_{1},\dots,T_{p}^{n_{p}}-z_{p})\quoi,
$$
of~$\othercal{V}$ on which $H$ acts, and the quotient of~$\othercal{V}''$
by $H$ is the desired extension.

The proof is then completed thanks to~\Ref{XII.4.6}.
The covering $\othercal{R}'$ comes from a finite covering $R'$
of~$R$; the restriction of~$R'$ to $X$ is a covering $X'$ of
$X$ such that we have $X^{\prime\an}\simeq \othercal{X}'$, and
by~\Ref{XII.3.1}(iii) $X'$ is an \'etale covering
of~$X$.

\begin{corollary}
\label{XII.5.2}
Let
% original p. 337
$X$ be a connected $\CC$-scheme locally of finite type,
$\varphi\colon X^{\an}{\to} X$ the canonical morphism, $x$ a point of
$X^{\an}$. Let $\pi_{1}(X^{\an},x)$ be the fundamental group of
the topological space $X^{\an}$ at the point $x$, and $\pi_{1}(X,\varphi(x))$
the fundamental group of the scheme $X$ at the point $\varphi(x)$
\textup{(V~\Ref{V.7})}. Then $\pi_{1}(X,\varphi(x))$ is canonically
isomorphic to the completion of~$\pi_{1}(X^{\an},x)$ for the
topology of subgroups of finite index.
\end{corollary}

Indeed, let $\cal{C}$ be the category of finite \'etale
coverings of~$X^{\an}$, $F$ the functor from~$\cal{C}$ to
$\Ens$ which, to every finite \'etale covering $\othercal{X}'$ of
$X^{\an}$, associates the set of points of~$\othercal{X}'$ over
$x$, and let $\widehat{\pi}_{1}(X^{\an},x)$ be the profinite group
associated with $\cal{C}$ and $F$ as explained in~V~\Ref{V.4}.
Since every finite \'etale covering of~$X^{\an}$ is a quotient of the
universal covering by a subgroup of finite index,
$\widehat{\pi}_{1}(X^{\an},x)$ is none other than the completion of
$\pi_{1}(X^{\an},x)$ for the topology of subgroups of finite index.
The corollary therefore follows from~\Ref{XII.5.1} and~V~\Ref{V.6.10}.

\begin{proposition}
\label{XII.5.3}
Let $\othercal{X}$ be a normal analytic space, $\othercal{Y}$ a
closed analytic subset such that
$\othercal{U}=\othercal{X}-\othercal{Y}$ is dense in $\othercal{X}$.
Then the functor which, to every normal finite
covering~\eqref{XII.5.0} $\othercal{X}'$ of~$\othercal{X}$, associates its restriction
to $\othercal{U}$ is fully faithful.
\end{proposition}

Let $\othercal{X}'$ and $\othercal{X}''$ be two normal finite coverings
of~$\othercal{X}$. We must show that the canonical map
$$
\Hom_{\othercal{X}}(\othercal{X}',\othercal{X}'') \to
\Hom_{\othercal{U}}(\othercal{X}'|\othercal{U},\othercal{X}''|\othercal{U})
$$
is bijective. Let $u$, $v$ be two $\othercal{X}$-morphisms from
$\othercal{X}'$ to $\othercal{X}''$ whose restrictions to $\othercal{U}$
are the same, and let us prove that $u=v$. The morphisms $u$ and
% original p. 338
$v$ coincide on the dense open set
$\othercal{U}\times_{\othercal{X}}\othercal{X}'$, hence on the underlying
topological spaces. By \hbox{\cite[\no 19~\S 4~cor\ptbl 5]{XII.4}} this
proves that we have $u=v$.

Now let $u$ be a $\othercal{U}$-morphism from~$\othercal{X}'|\othercal{U}$
to $\othercal{X}''|\othercal{U}$, and let us show that it extends to
the whole of~$\othercal{X}'$. We may assume $\othercal{X}'$ regular.
Indeed, $\othercal{X}'$ being normal, one can find an open set
$\othercal{V}$ of~$\othercal{X}$ whose complement is an analytic
subset of codimension $\geq 2$, such that
$\othercal{X}'\times_{\othercal{X}}\othercal{V}=\othercal{V}'$ is regular.
Let $\othercal{V}''=\othercal{X}''\times_{\othercal{X}}\othercal{V}$ and suppose the
proposition proved for $\othercal{V}$. Consider the commutative
diagram
$$
\xymatrix{ \othercal{V}' \ar[dd]_{g'}\ar[rd]\ar[rrr]^{i'} && &\othercal{X}'
\ar[dd]^{\hbox{\raise10mm\hbox{$f'$}}} \cr &\othercal{V}'' \ar[ld]_{g''\mkern-10mu}\ar[rrr]^-{i''~~} & &
&\othercal{X}'' \ar[ld]_{f''\mkern-12mu} \cr \othercal{V} \ar[rrr]^{i} && &\othercal{X} \cr
}
$$
To $u$ is associated a morphism of
$\cal{O}_{\othercal{V}}$-Algebras $g''_*\cal{O}_{\othercal{V}''}\to
g'_*\cal{O}_{\othercal{V}'}$, from which one deduces a morphism
$$
i_*g''_*\cal{O}_{\othercal{V}''} \to
i_*g'_*\cal{O}_{\othercal{V}'}\quoi.
$$
Taking into account the isomorphisms
$i'_*\cal{O}_{\othercal{V}'}\simeq\cal{O}_{\othercal{X}'}$,
$i''_*\cal{O}_{\othercal{V}''}\simeq\cal{O}_{\othercal{X}''}$
\cite[\no 3~prop\ptbl 4]{XII.11}, one deduces from it a morphism of
$\cal{O}_{\othercal{X}}$-Algebras
$$
f''_*\cal{O}_{\othercal{X}''} \to
f'_*\cal{O}_{\othercal{X}'}\quoi,
$$
whence the desired morphism $\othercal{X}'\to\othercal{X}''$.

From
% original p. 339
now on we assume $\othercal{X}'$ regular. Let
$\othercal{U}'=\othercal{U}\times_{\othercal{X}}\othercal{X}'$,
$\othercal{Y}'=\othercal{X}'-\othercal{U}'$. We regard $\othercal{Y}'$ as a
reduced analytic subspace of~$\othercal{X}'$; if $\othercal{Y}'_{1}$
is the singular closed set of~$\othercal{Y}'$, we have $\dim \othercal{Y}'_{1} <
\dim \othercal{Y}'$~\cite[\no 20~D\ptbl Th\ptbl 3]{XII.4}. One therefore sees by induction
on the dimension of~$\othercal{Y}'$ that we may assume $\othercal{Y}'$
smooth. Since it suffices to extend $u$ to an open neighborhood of
each point of~$\othercal{Y}'$, we may assume, by the implicit function
theorem, that $\othercal{X}'$ is a ball of an affine
space $\othercal{E}^{n}$ and $\othercal{Y}'$ the closed set defined by
the vanishing of the first $p$ coordinate functions
$z_{1},\dots,z_{p}$, with $0\leq p\leq n$.

To $u$ we associate a section $s$ of
$p\colon\othercal{X}'\times_{\othercal{X}}\othercal{X}''\to\othercal{X}'$ over
$\othercal{U}'$; after restricting $\othercal{X}'$, we may assume
$p_*(\cal{O}_{\othercal{X}'\times_{\othercal{X}}\othercal{X}''})$
generated by elements $x_{1},\dots,x_{q}$ of
$\Gamma(\othercal{X}',
p_*\cal{O}_{\othercal{X}'\times_{\othercal{X}}\othercal{X}''})$; let
$u_{1},\dots,u_{q}\in\Gamma(\othercal{U}',\cal{O}_{\othercal{X}'})$ be the images
under $s$ of~${x_{1}}|\othercal{U}',\dots,{x_{q}}|\othercal{U}'$. To say
that $s$ extends to $\othercal{X}'$ amounts to saying that the
$u_{1},\dots,u_{q}$ extend to sections of
$\Gamma(\othercal{X}',\cal{O}_{\othercal{X}'})$. But, since $f$ is finite,
each $u_{i}$ is a Laurent series in $z_{1},\dots,z_{p}$,
with coefficients power series in
$z_{p{+}1},\dots,z_{n}$, which satisfy relations of
integral dependence. It follows that $u_{i}$ is
bounded, hence is a power series in $z_{1},\dots,z_{n}$,
and consequently extends to $\othercal{X}'$.

One may ask whether the functor introduced in~\Ref{XII.5.3} is
an equivalence of categories. An answer to
this question is provided by the theorem of \textsc{Grauert}-\textsc{Remmert} \cite{XII.6},
of which we give a proof below using
resolution of singularities. One could also have used the
theorem of \textsc{Grauert}-\textsc{Remmert} to prove~\Ref{XII.5.1};
this is what was done before~\cite{XII.8} was available.

\begin{theorem}[\textsc{Grauert}-\textsc{Remmert} theorem]
\label{XII.5.4}
Let
% original p. 340
$\othercal{X}$ be a normal analytic space, $\othercal{Y}$ a
closed analytic subset such that
$\othercal{U}=\othercal{X}-\othercal{Y}$ is dense in~$\othercal{X}$. Let
$\othercal{U}'$ be a normal finite covering of~$\othercal{U}$; assume
that there exists a nowhere dense closed analytic subset $\othercal{S}$
of~$\othercal{X}$ such that the restriction of~$\othercal{U}'$ to
$\othercal{U}-\othercal{U}\cap\othercal{S}$ is \'etale. Then there exists a
normal finite covering $\othercal{X}'$ of~$\othercal{X}$ which extends
$\othercal{U}'$, and $\othercal{X}'$ is unique up to isomorphism.
\end{theorem}

Uniqueness follows from~\Ref{XII.5.3}. The problem of
extending $\othercal{U}'$ is therefore local on~$\othercal{X}$. We may assume
$\othercal{U}$ regular and $\othercal{U}'$ \'etale over~$\othercal{U}$.
Indeed the set of regular points of~$\othercal{U}$ is an open set
$\othercal{V}$ dense in~$\othercal{X}$ whose complement is an
analytic subset \hbox{\cite[\no20~D th\ptbl 2]{XII.4}}, and it suffices to replace
$\othercal{U}$ by the open set $\othercal{V}-\othercal{V}\cap\othercal{S}$.

Let $y$ be a point of~$\othercal{X}-\othercal{U}$ and let us show that
$\othercal{U}'$ can be extended to a neighborhood of~$y$. After
restricting $\othercal{X}$ to an open neighborhood of~$y$, it
follows from the resolution of singularities theorem
\cite{XII.8} that one can find a regular analytic space
$\othercal{X}_1$ and a projective morphism $f\colon\othercal{X}_1\to\othercal{X}$
inducing by restriction to~$\othercal{U}$ an isomorphism
$\othercal{U}_1=f^{-1}(\othercal{U})\simeq\othercal{U}$, such that $\othercal{U}_1$ is
the complement in~$\othercal{X}_1$ of a divisor with
normal crossings. Let us show that $\othercal{U}'$ extends to a
normal finite covering of~$\othercal{X}_1$. Since the question is
local on~$\othercal{X}_1$, we may assume that $\othercal{X}_1$ is a
ball of an affine space $\othercal{E}^n$ and that $\othercal{X}_1-\othercal{U}_1$
is defined by the vanishing of the first $p$ coordinate
functions
$z_1,\dots,z_p$,
with $0\leq p\leq n$. The \'etale covering $\othercal{U}'$
of~$\othercal{U}_1$ is a quotient of a covering of the form
$$
\othercal{U}_2=\othercal{U}_1[T_1,\dots,T_p]/\left(T_1^{n_1}-z_1,\dots,
T_p^{n_p}-z_p\right)
$$
by a subgroup $H$ of the Galois group of~$\othercal{U}_2$. The
covering $\othercal{U}_2$ extends
% original p. 341
to the covering
$$
\othercal{X}_2=\othercal{X}_1[T_1,\dots,T_p]/\left(T_1^{n_1}-z_1,\dots,
T_p^{n_p}-z_p\right)
$$
of~$\othercal{X}_1$ on which $H$ acts, and $\othercal{X}_2/H$ extends
$\othercal{U}'$ to~$\othercal{X}_1$.

Denote by $\othercal{X}_1'$ the normal finite covering of~$\othercal{X}_1$ which
extends $\othercal{U}'$, and by $\othercal{F}_1$ the coherent
$\cal{O}_{\othercal{X}_1}$-Algebra defined
by~$\othercal{F}_1$. By the finiteness theorem of
\textsc{Grauert}-\textsc{Remmert} \hbox{\cite[\no15 th\ptbl 1.1]{XII.4}}, $f_*\othercal{F}_1$ is a
coherent $\cal{O}_{\othercal{X}}$-Algebra. There therefore corresponds to it
a finite covering $\othercal{X}'$ of~$\othercal{X}$, which is moreover
normal since $\othercal{X}_1'$ is, and $\othercal{X}'$ is the desired
extension of~$\othercal{U}'$.

\begin{remark}
\label{XII.5.5}
In statement~\Ref{XII.5.4}, one cannot drop
the hypothesis on the locus of points where the morphism
$\othercal{U}'\to\othercal{U}$ is not \'etale. For example, let
$\othercal{X}$ be the unit disk of the complex plane, $\othercal{U}$ the
complement of the origin in~$\othercal{X}$,
$\othercal{U}'=\othercal{U}[T]/(T^2-\sin 1/z)$, where~$z$ is the
coordinate function on~$\othercal{X}$. Then $\othercal{U}'$ is a normal
finite covering of~$\othercal{U}$ which does not extend
to~$\othercal{X}$. Indeed, suppose that $\othercal{U}'$ extends to a
finite covering $\othercal{X}'$ of~$\othercal{X}$; the locus of points
of~$\othercal{X}$ where the morphism $\othercal{X}'\to\othercal{X}$ is not
\'etale is then a closed analytic set which contains all the
points $z$ such that $\sin 1/z=0$, which is absurd.

One can, however, drop the hypothesis on the singular locus of the
morphism $\othercal{U}'\to\othercal{U}$ when we have
$\codim(\othercal{X}-\othercal{U},\othercal{X})\geq~2$. Indeed, we may assume
$\othercal{U}$ regular. The locus of points of~$\othercal{U}$ where
$\othercal{U}'\to\othercal{U}$ is not \'etale is a divisor
of~$\othercal{U}$, and it follows from the theorem of \textsc{Remmert}-\textsc{Stein}
\cite[th\ptbl 3]{XII.9} that it is
% original p. 342
the trace on~$\othercal{U}$ of a divisor of~$\othercal{X}$. Now, in this case,
if $\othercal{A}$ is a coherent $\cal{O}_{\othercal{U}}$-Algebra
such that $\othercal{U}'=\Spec \an(\othercal{A})$, and if $i:\othercal{U}\to\othercal{X}$
is the canonical morphism, it suffices to take
$\othercal{X}'=\Spec\an(i_*\othercal{A})$; indeed, we know that $i_*\othercal{A}$
is coherent \cite[\no 1 th\ptbl 1]{XII.11}.
\end{remark}

\begin{remarkMR}
\label{XII.5.6}
There exist nontrivial groups $G$ of finite presentation
which possess no subgroups of finite index
distinct from~$G$, for example the group of G\ptbl Higman (\cf J-P\ptbl Serre, Arbres et amalgames, Ast\'erisque \No 46, prop\ptbl 6, chap\ptbl I, \S 1).
Consequently, if such a group is realized as the topological fundamental
group of a scheme $V$ over~$\CC$, say smooth and
projective, the topological space $V_{\mathrm{top}}$ underlying $V$ is not
simply connected, but $V$ is algebraically simply connected.
At present, no such $V$ is known. Let us point out, however,
that D\ptbl Toledo has constructed smooth projective schemes $V$ over~$\CC$ whose topological fundamental group is not separated
for the topology of subgroups of finite index (the natural morphism
$\pi_1(V_{\mathrm{top}})\to\pi_1(V_{\alg})$ is not injective). [D\ptbl Toledo, Projective varieties with non-residually finite fundamental group, Publ.\ Math.\ IHES \textbf{77} (1993), p\ptbl103--119].
\end{remarkMR}

\begin{thebibliography}{10}{XII.6}
% original p. 343

\bibitem[1]{XII.1} N\ptbl \textsc{Bourbaki},
Topologie G\'en\'erale, Hermann, Paris (1960).

\bibitem[2]{XII.2} N\ptbl \textsc{Bourbaki},
Alg\`ebre Commutative, Hermann, Paris (1961).

\bibitem[3]{XII.3} H\ptbl \textsc{Cartan}. S\'eminaire E.N.S., Paris (1956-57).

\bibitem[4]{XII.4} H\ptbl \textsc{Cartan}. S\'eminaire E.N.S., Paris (1960-61).

\bibitem[5]{XII.5} R\ptbl \textsc{Godement}, Th\'eorie des Faisceaux, Hermann, Paris
(1958).

\bibitem[6]{XII.6} H\ptbl \textsc{Grauert} and R\ptbl \textsc{Remmert},
Komplexe R\"aume, Math.\ Ann. \textbf{136} (1958), p\ptbl 245--318.

\bibitem[7]{XII.7} M\ptbl \textsc{Hakim}, Sch\'emas relatifs, thesis, Paris (1967).

\bibitem[8]{XII.8} H\ptbl \textsc{Hironaka},
Resolution of singularities of an algebraic variety over a field of characteristic zero, Ann.\ of Math. \textbf{39} (1964), p\ptbl 109--236.

\bibitem[9]{XII.9} R\ptbl \textsc{Remmert} and K\ptbl \textsc{Stein},
Ueber die wesentlichen Singularit\"aten analytischer Mengen,
Math.\ Ann. \textbf{126} (1953), p\ptbl 263--306.

\bibitem[10]{XII.10}
J-P\ptbl\textsc{Serre},
G\'eom\'etrie alg\'ebrique et g\'eom\'etrie analytique, Ann.\ Inst.\ Fourier (Grenoble) \textbf{6} (1956), p\ptbl 1--42.

\bibitem[11]{XII.11}
J-P\ptbl\textsc{Serre},
Prolongement de faisceaux analytiques
coh\'erents, Ann.\ Inst.\ Fourier (Grenoble) \textbf{16} (1966), p\ptbl 363--374.

\end{thebibliography}
